\section{Introducción}

\subsection{Resumen}
Este documento presenta el análisis, diseño, implementación y verificación del \textbf{analizador sintáctico} para kernels \textbf{Triton GPU} escritos en archivos \texttt{.tri}. La fase sintáctica es la segunda etapa del compilador del reto académico: consume la secuencia de tokens producida por el analizador léxico (Flex) del Equipo~11, construye una estructura conforme a una gramática libre de contexto (CFG) implementada con \textbf{yacc/Bison}, actualiza tablas de símbolos y reporta errores sintácticos.

El entregable incluye: (i) gramática documentada (Apéndice~\ref{ap:cfg-completa}, CFG completa); (ii) tablas de símbolos; (iii) implementación \texttt{triton\_parser}; (iv) mensajes de error; (v) \textbf{resultados experimentales} con énfasis en la Sección~\ref{sec:resultados}: \textbf{100\%} suite unitaria (12/12), \textbf{100\%} suites JSONL (200/200), benchmark TRITONHEAL (58/70) y gráfica de los 70 kernels.

\subsection{Notación}
\begin{itemize}
  \item \textbf{CFG}: gramática libre de contexto en notación BNF; producciones $A \rightarrow \alpha$.
  \item \textbf{Token ID}: identificador numérico del Cuadro~1 del reporte léxico (100--999).
  \item \textbf{INDENT/DEDENT} (910/911): tokens sintéticos derivados de \texttt{WS\_COUNT} (900) para modelar indentación Python.
  \item \textbf{PARSE\_OK}: análisis sintáctico exitoso con impresión de tabla de símbolos del parser.
\end{itemize}

\subsection{Justificación de herramientas}
Se utiliza \textbf{C} con \textbf{Flex} y \textbf{yacc/Bison}, conforme a las restricciones del curso~[1]. No se emplean librerías CFG de alto nivel ni \texttt{ast} de Python para el parser formal. El modelo de desarrollo sigue el proceso de análisis, diseño, implementación y verificación definido en el IEEE-830~[4], con los cinco entregables del Entregable~B como unidades de documentación.

\subsection{Trazabilidad a la fase léxica}
Los IDs de token se preservan en \texttt{\%token} de \texttt{parser.y} (p.\,ej.\ \texttt{TOKEN\_IDENTIFIER = 100}). El archivo \texttt{lexico\_equipo\_11\_scanner.l} es la base del scanner integrado; el modo \texttt{-t} reproduce la salida del reporte léxico (TC-01--TC-04).
